\documentclass[11pt,letterpaper]{article}

% ------ Configuración de Idioma y Codificación ------
\usepackage[utf8]{inputenc}
\usepackage[spanish, es-tabla]{babel} % es-tabla cambia "Cuadro" por "Tabla"

% ------ Paquetes de Color y Estilos Visuales ------
\usepackage[table]{xcolor} % Debe ir arriba para evitar conflictos con tablas
\usepackage{tcolorbox}
\usepackage{epsfig,graphicx}
\usepackage{tikz}

% ------ Configuración de Geometría y Páginas ------
\usepackage[left=2cm,right=2cm,top=1.8cm,bottom=2.3cm]{geometry}
\usepackage{fancyhdr}
\usepackage{lastpage}
\pagestyle{fancy}
\fancyhf{}
\fancyhead[R]{\ifnum\value{page}>1 \textbf{Estadística Bayesiana}\fi}
\lfoot{\includegraphics[height=1.0cm]{ciencias.png}}
\rfoot{\textit{Página \thepage}}

% ------ Paquetes Matemáticos Básicos y Teoremas ------
\usepackage{amsmath}
\usepackage{amssymb}
\usepackage{amsthm}

\newtheorem{theorem}{Theorem}[section]
\newtheorem{corollary}{Corollary}[theorem]
\newtheorem{lemma}[theorem]{Lemma}

% ------ Tablas Avanzadas ------
\usepackage{tabularx}
\usepackage{colortbl}
\usepackage{array, multirow, multicol}

% ------ Formateo de Código (Listings) ------
\usepackage{listings}

% Definición de Colores para el Código
\definecolor{codegreen}{rgb}{0,0.6,0}
\definecolor{codegray}{rgb}{0.5,0.5,0.5}
\definecolor{codepurple}{rgb}{0.58,0,0.82}
\definecolor{backcolour}{rgb}{0.95,0.95,0.92}
\definecolor{maroon}{HTML}{7B1E3A}
\definecolor{AzulUNAM}{RGB}{12, 35, 64}     % Azul marino institucional
\definecolor{OroUNAM}{RGB}{191, 144, 0}     % Oro institucional (el que ya tenías)
\definecolor{backcolour}{rgb}{0.95,0.95,0.92}

% Configuración ÚNICA de mystyle
\lstdefinestyle{mystyle}{
    backgroundcolor=\color{backcolour},
    commentstyle=\color{codegreen},
    keywordstyle=\color{magenta},
    numberstyle=\tiny\color{codegray},
    stringstyle=\color{codepurple},
    basicstyle=\ttfamily\footnotesize,
    breakatwhitespace=false,
    breaklines=true,
    captionpos=b,
    keepspaces=true,
    numbers=left,
    numbersep=5pt,
    showspaces=false,
    showstringspaces=false,
    showtabs=false,
    tabsize=2
}
\lstset{style=mystyle}

% ------ Algoritmos y Pseudocódigo ------
\usepackage{algorithm}
\usepackage{algpseudocode}

% ------ Hipervínculos y Botón de Copiado (Al Final) ------
\usepackage{url}
\usepackage{hyperref}
\usepackage{clipboard} 
% ------ Atajos Matemáticos (Macros) ------
\definecolor{B}{HTML}{FFFFFF}
\definecolor{G}{HTML}{5e5e5e}
\definecolor{R2}{HTML}{d53d40}
\definecolor{A2}{HTML}{034190}
\definecolor{V2}{HTML}{7faa50}

\newcommand{\R}{\mathbb{R}}
\newcommand{\C}{\mathcal{C}}
\newcommand{\N}{\mathbb{N}}
\newcommand{\Z}{\mathbb{Z}}
\newcommand{\Q}{\mathbb{Q}}

% ------ Configuración de Listas Enumeradas ------
\renewcommand{\theenumi}{\Roman{enumi}}
\renewcommand{\labelenumi}{\theenumi.}

\begin{document}

%---Evironment--%
	\makeatletter
        \renewenvironment{proof}[1][\proofname]{\par
            \pushQED{\qed}%
            \normalfont \topsep6\p@\@plus6\p@\relax
            \trivlist
            \item\relax
            {\itshape
            #1\@addpunct{.}}\par\vspace{\baselineskip}\ignorespaces
            }{%
            \popQED\endtrivlist\@endpefalse
            }
        \newenvironment{Solucion}[1][Soluci\'on]{\par
            \pushQED{\qed}%
            \normalfont \topsep6\p@\@plus6\p@\relax
            \trivlist
            \item\relax
            {\itshape
            #1\@addpunct{.}}\par\vspace{\baselineskip}\ignorespaces
            }{%
            \popQED\endtrivlist\@endpefalse
            }
            
    \makeatother

%------ Encabezado -------- %

\begin{center}
    \begin{minipage}{3cm}
    	\begin{center}
    		\includegraphics[height=3.5cm]{unam.jpg}
    	\end{center}
    \end{minipage}\hfill
    \begin{minipage}{10cm}
    	\begin{center}
    	\textbf{\large Universidad Nacional Autónoma de México}\\[0.1cm]
        \textbf{Facultad de Ciencias}\\[0.1cm]
        \textbf{\'Estadistica Bayesiana}\\[0.1cm]
        Tareita 8,9,10 \\[0.1cm]
         El\'ias L\'opez Rivera\\[0.1cm]
        \texttt{ elias.lopezr,@ciencias.unam.mx }\\[0.1cm]
        Fecha:\,\,24/09/2026
    	\end{center}
    \end{minipage}\hfill
    \begin{minipage}{3cm}
    	\begin{center}
    		\includegraphics[height=3.5cm]{ciencias.png}
    	\end{center}
    \end{minipage}
\end{center}
\rule{17cm}{0.1mm}
%------ Fin de encabezado -------- %
\begin{tcolorbox}[
    title = \textcolor{white}{\textbf{Problema 1}},
    colbacktitle=AzulUNAM!95!black,         % Fondo del título en azul marino profundo
    coltitle=white,                         % Texto del título en blanco
    colback=OroUNAM!7!white,                % Fondo del cuerpo con un ligero toque de oro
    colframe=OroUNAM!90!black,              % Borde del recuadro en oro institucional
    fontupper=\small
]
\textit{En un ejercicio de dos valores esperados de una Gamma, comparaci\'on de las funciones $\phi$ }
\end{tcolorbox}
Tomamos las dos predictivas encontradas en clase para el caso del doctor $P_c(\theta|t_m^c)=ga(\theta|21.5,182.125)$ y $P_t(\theta|t_m^t)=Ga(\theta|21,326)$, recordamos que la
cantidad de interes es $\pi=\frac{1}{\theta}$, simularemos muestras de ambas distribuciones y con ellas de la distribcui\'on final de $\phi_t-\phi_c$ 
 \begin{lstlisting}[style=mystyle, language=Python, caption={Simulacion Integral y verosimilitud}]
import numpy as np

# Configuracion de la simulacion masiva
N_sim = 500000
np.random.seed(42)

# ==============================================================================
# 1. GRUPO CONTROL: P_c(theta | t_m^c) = Ga(shape = 21.5, rate = 182.125)
# ==============================================================================
# En numpy, el parametro 'scale' es la inversa de la tasa (scale = 1 / rate)
theta_c = np.random.gamma(shape=21.5, scale=1/182.125, size=N_sim)
phi_c = 1 / theta_c

# ==============================================================================
# 2. GRUPO TRATAMIENTO: P_t(theta | t_m^t) = Ga(shape = 21, rate = 326)
# ==============================================================================
theta_t = np.random.gamma(shape=21, scale=1/326, size=N_sim)
phi_t = 1 / theta_t

# ==============================================================================
# 3. CANTIDAD DE INTERES: Diferencia de los parametros transformados
# ==============================================================================
diferencia_phi = phi_t - phi_c

# ==============================================================================
# 4. IMPRESION DE ESTADISTICOS POSTERIORES
# ==============================================================================
print(f"----- ANALISIS POSTERIOR DEL CASO DEL DOCTOR -----")
print(f"Media a posteriori de phi_c (Control)    : {np.mean(phi_c):.4f}")
print(f"Media a posteriori de phi_t (Tratamiento): {np.mean(phi_t):.4f}")
print(f"Media a posteriori de la diferencia       : {np.mean(diferencia_phi):.4f}")
print(f"Intervalo de credibilidad del 95%         : [{np.percentile(diferencia_phi, 2.5):.4f}, {np.percentile(diferencia_phi, 97.5):.4f}]")
\end{lstlisting}
\begin{tcolorbox}[
    title = \textcolor{white}{\textbf{Problema 2}},
    colbacktitle=AzulUNAM!95!black,         % Fondo del título en azul marino profundo
    coltitle=white,                         % Texto del título en blanco
    colback=OroUNAM!7!white,                % Fondo del cuerpo con un ligero toque de oro
    colframe=OroUNAM!90!black,              % Borde del recuadro en oro institucional
    fontupper=\small]
\textit{ Al establecer una distribución Normal como la inicial de una dsitribución Cauchy, obtener el valor de la integral considerando al producto de la dsitribuciones como el valor esperado de una función de la normal o como el valor esperado de una función de la Cauchy. Utilizar la ley de los grande números para simular y aproximar el valor esperado con la media aritmética.}
\end{tcolorbox}
\begin{proof}
    Supongamos $s_m=\{x_1,\cdots,x_m\}\sim Ca(x|\theta,1)$, y $\theta\sim N(\theta|0,1)$, por tanto:\,\\
    \begin{equation*}
        p(\theta|s_m)\varpropto e^{-\frac{\theta^2}{2}}\left(\pi\displaystyle\prod_{j=1}^m(1+(x_j-\theta)^2)\right)^{-1} 
    \end{equation*}
    por tanto $I=\int_{-\infty}^{\infty}\frac{e^{-\frac{\theta^2}{2}}}{\sqrt{2\,\pi}}\left(\pi\displaystyle\prod_{j=1}^m(1+(x_j-\theta)^2)\right)^{-1}d\theta$ es su constante de normalizaci\'on, claramente converge pues:\,\\
    \begin{align*}
        \left(\displaystyle\prod_{j=1}^m\pi(1+(x_j-\theta)^2)\right)^{-1}\leq 1\implies I\leq \int_{-\infty}^{\infty}e^{-\frac{\theta^2}{2}}d\theta<\infty
    \end{align*} 
    $I$ no tiene soluci\'on an\'alitica por tanto la paroximaremos de dor formas diferentes:\,\\
    \begin{enumerate}
        \item $f(\theta)=\frac{e^{-\frac{\theta^2}{2}}}{\sqrt{2\pi}}$, por tanto usando la ley del estadistico ciego:\,\\
        \begin{equation*}
            I=\int_{-\infty}^{\infty}\left(\pi\displaystyle\prod_{j=1}^m(1+(x_j-\theta)^2)\right)^{-1}f(\theta)=E_{\theta}\left[\left(\displaystyle\prod_{j=1}^m\pi(1+(x_j-\theta)^2)\right)^{-1}\right]=E_{\theta}[g(\theta)]
        \end{equation*}\,\\
        Como $I$ existe entonces $E_{\theta}[g(\theta)]<\infty$ podemos aplicar la ley de grandes n\'umeros, si tomamos $\{\bar{\theta}_i\}_{i\in \N}\sim g(\theta) $ identicamente distribuidas e independientes con media constante tenemos que:\,\\
        \begin{equation*}
            P\left(\lim_{n\to\infty}\frac{1}{n}\displaystyle\sum_{i=1}^n \bar{\theta_i}=E_{\theta}(g(\theta))\right)=1
        \end{equation*}
        por tanto para $N$ suficientemente grande:
        \begin{equation*}
            I\approx \displaystyle\frac{1}{N}\sum_{i=1}^{N}\left(\pi\displaystyle\prod_{j=1}^m(1+(x_j-\theta_i)^2)\right)^{-1}
        \end{equation*}
        
        \item $f_k(\theta)=\frac{1}{\pi(1+(x_k-\theta)^2)}$, de nuevo usando la ley del estadistico ciego:\,\\
        \begin{align*}
            I&=\int_{-\infty}^{\infty}f_k(\theta)e^{-\frac{\theta^2}{2}}\left(\displaystyle\pi\prod_{j\neq k\,\,j=1}^m(1+(x_j-\theta^2))\right)^{-1}d\theta= E_{Ca(x_k,1)}\left[\left(e^{-\frac{\theta^2}{2}}\displaystyle\prod_{j\neq k\,\,j=1}^m\pi(1+(x_j-\theta^2))\right)^{-1}\right]\,\\
            \,\\&=E_{Ca(x_k,1)}[h(\theta)]
        \end{align*}\,\\
        Como $I$ existe entonces $E_{Ca(x_k,1)}[h(\theta)]<\infty$ podemos aplicar la ley de grandes n\'umeros, si tomamos $\{\bar{\theta}_1\}_{i\in \N}\sim h(\theta)$ identicamente distribuidas e independientes con media constante tenemos que:\,\\
        \,\\
        \begin{equation*}
            P\left(\lim_{n\to\infty}\frac{1}{n}\displaystyle \sum_{i=1}^n\bar{\theta_i}=E_{ca(\theta_k,1)}(h(\theta))\right)=1
        \end{equation*}\,\\
        \,\\
        Por tanto para $N$ suficientemente grande:\,\\
        \begin{equation*}
            I\approx\frac{1}{N}\displaystyle \sum_{i=1}^Ne^{-\frac{\theta_i^2}{2}}\left(\displaystyle\prod_{j\neq k\,\,j=1}^m\pi(1+(x_j-\theta_i^2))\right)^{-1}
        \end{equation*}
    \end{enumerate}
    \begin{lstlisting}[style=mystyle, language=Python, caption={Simulacion Integral y verosimilitud}]
import numpy as np

np.random.seed(42)
tam_muestra = 10   # Tamano de la muestra de datos observados (m en tu planteamiento)
M_sim = 100000     # Numero de replicas de Monte Carlo para aplicar la LGN
centro_real = 2

# Generamos la muestra observada real que distribuye Cauchy(centro_real, 1)
X_datos = centro_real + np.random.standard_cauchy(size=tam_muestra)
X_fila = X_datos[np.newaxis, :]  # Forma: (1, tam_muestra)

# ==============================================================================
# ENFOQUE A: DISTRIBUCION INICIAL = NORMAL ESTANDAR
# Integral = E_{Normal}[ Producto_{i=1}^m Cauchy(X_i | theta) ]
# ==============================================================================
thetas_A = np.random.normal(loc=0.0, scale=1.0, size=M_sim)
thetas_col_A = thetas_A[:, np.newaxis]  # Forma: (M_sim, 1)

# Log-verosimilitud conjunta de los datos (sin incluir la Normal porque es la FDP base)
log_Mat_cauchy_A = - np.log(np.pi) - np.log(1 + (X_fila - thetas_col_A)**2)
log_objetivo_A = np.sum(log_Mat_cauchy_A, axis=1)

# Estabilizacion numerica Log-Sum-Exp para el Enfoque A
max_log_A = np.max(log_objetivo_A)
evidencia_A_estable = np.exp(log_objetivo_A - max_log_A)
constante_A = np.mean(evidencia_A_estable) * np.exp(max_log_A)


# ==============================================================================
# ENFOQUE B: DISTRIBUCION INICIAL = CAUCHY COMPUESTA (centrada en X_datos)
# Integral = E_{Cauchy(X_0)}[ Normal(theta) * Producto_{i=2}^m Cauchy(X_i | theta) ]
# ==============================================================================
x_base = X_datos[0]
thetas_B = x_base + np.random.standard_cauchy(size=M_sim)
thetas_col_B = thetas_B[:, np.newaxis]  # Forma: (M_sim, 1)

# 1. Log-densidad del Prior Normal
log_prior_normal = - 0.5 * np.log(2 * np.pi) - (thetas_B**2 / 2)

# 2. Log-verosimilitud de las restantes m-1 observaciones de Cauchy (X_datos[1:] en adelante)
X_restantes_fila = X_datos[1:][np.newaxis, :] # Forma: (1, tam_muestra - 1)
log_Mat_cauchy_B = - np.log(np.pi) - np.log(1 + (X_restantes_fila - thetas_col_B)**2)

# Funcion objetivo conjunta para el enfoque B
log_objetivo_B = log_prior_normal + np.sum(log_Mat_cauchy_B, axis=1)

# Estabilizacion numerica Log-Sum-Exp para el Enfoque B
max_log_B = np.max(log_objetivo_B)
evidencia_B_estable = np.exp(log_objetivo_B - max_log_B)
constante_B = np.mean(evidencia_B_estable) * np.exp(max_log_B)


# ==============================================================================
# IMPRESION DE RESULTADOS
# ==============================================================================
print(f"----- CONSTANTE DE NORMALIZACION DE LA PREDICTIVA -----")
print(f"Tamano de muestra de datos: {tam_muestra}")
print(f"Replicas de Monte Carlo: {M_sim}\n")
print(f"Enfoque A (Normal Inicial)  : {constante_A:.18e}")
print(f"Enfoque B (Cauchy Inicial)  : {constante_B:.18e}")
\end{lstlisting}
\end{proof}

\begin{tcolorbox}[
    title = \textcolor{white}{\textbf{Problema 3}},
    colbacktitle=AzulUNAM!95!black,         % Fondo del título en azul marino profundo
    coltitle=white,                         % Texto del título en blanco
    colback=OroUNAM!7!white,                % Fondo del cuerpo con un ligero toque de oro
    colframe=OroUNAM!90!black,              % Borde del recuadro en oro institucional
    fontupper=\small]
\textit{Obtener conjugada de una uniforme $S_m\sim Unif(x|0,\theta )$.
}
\end{tcolorbox}
\begin{proof}
    Por lo visto en clase la conjugada de $S_m\sim Unif(x|0,\theta)$ es $pa(0,\theta|\alpha,\beta_2)$ es decir:\,\\
    \begin{equation*}
        pa(\theta|\alpha,\beta)=\alpha\beta^{\alpha}\theta^{-(\alpha+1)}I_{(\beta,\infty)}(\theta)
    \end{equation*}
    Es decir :\,\\
    \begin{align*}
        p(\theta|S_m)&\varpropto \theta^{-(\alpha+m+1)}I_{(x_{(m)},\infty)}(\theta)I_{(x_{(,\infty)}(\theta))}\,\\
        \,\\
        &=pa(\theta|\alpha+m,max\{\beta,x_{(m)}\})
    \end{align*}\,\\
    Donde $x_{(m)}$ es el m\'aximo de la muestra
\end{proof}
\end{document}